% Introduction and Overview  (source §1; added: reading guide, dependency
% diagram as a labeled figure, main theorem stated up front)
\chapter{Introduction and Overview}
\label{ch:intro}

\section{A Construction That Should Not Work}
\label{sec:intro-construction}

Take a smooth cubic curve in the plane, for instance
\[
  E:\quad y^2 = 4x^3 - g_2 x - g_3,
\]
and the following recipe for ``adding'' two of its points. Given points
\(P\) and \(Q\) on \(E\), draw the straight line through them. A line meets a
cubic in three points, so it strikes \(E\) at one further point \(R\). Reflect
\(R\) across the \(x\)-axis to obtain a point we call \(P \oplus Q\). To add a
point to itself, use the tangent line at \(P\) in place of the chord. This is
the \emph{chord--tangent law}.

Nothing about this recipe looks like it should produce a group. Reflecting a
third intersection point is a geometric gesture, not an algebraic operation;
there is no visible reason for it to be associative, and no visible reason for
inverses to behave. Yet the chord--tangent law \emph{is} an abelian group law,
with associativity holding exactly. The standard elementary proofs of
associativity are long case analyses or invocations of the Cayley--Bacharach
theorem, and they leave the impression that the group structure is a lucky
accident of plane geometry.

It is not an accident. These notes explain the following single fact, which
makes every group axiom transparent at once.

\begin{quote}
The chord--tangent law on a smooth plane cubic is not an arbitrary geometric
trick. It is exactly the abelian group law on the complex torus
\(T = \C/\Lambda\), transported to the cubic by the Weierstrass functions
\(\wp\) and \(\wp'\).
\end{quote}

The slogan is
\[
  \boxed{
    \text{addition on } \C/\Lambda
    \quad\Longleftrightarrow\quad
    \text{chord--tangent addition on } E
  }
\]
where, in homogeneous coordinates \([X:Y:Z]\) on \(\CP^2\),
\[
  E:\quad Y^2 Z = 4X^3 - g_2 X Z^2 - g_3 Z^3 .
\]
On the torus side, addition is nothing more than addition of complex numbers
modulo a lattice \(\Lambda\); it is associative and commutative for the most
elementary of reasons. The entire mystery of the chord--tangent law dissolves
once we exhibit a dictionary between the two pictures.

\section{The Answer Stated First}
\label{sec:intro-answer}

We state the destination now, before any of the machinery is built, so that
every later chapter can be read as a step toward it. The terms appearing here
(\emph{lattice}, \emph{Weierstrass function}, \emph{divisor}, \emph{point at
infinity}) are defined from scratch in the chapters that follow; forward
references point to the precise statements.

\begin{theorem}[Main theorem]
\label{thm:main}
Let \(\Lambda \subset \C\) be a lattice and let \(\wp\) be the Weierstrass
function of \(\Lambda\), with derivative \(\wp'\). Set
\[
  E:\quad Y^2 Z = 4X^3 - g_2 X Z^2 - g_3 Z^3 \subset \CP^2 ,
\]
with \(g_2, g_3\) the invariants of \(\Lambda\). Then the map
\[
  \Phi : \C/\Lambda \longrightarrow E,
  \qquad
  \Phi(z) =
  \begin{cases}
    [\wp(z) : \wp'(z) : 1], & z \notin \Lambda,\\[2pt]
    [0 : 1 : 0] =: O,       & z \in \Lambda,
  \end{cases}
\]
is an isomorphism of compact Riemann surfaces, and it is an isomorphism of
groups when \(E\) carries the chord--tangent law and \(\C/\Lambda\) carries
addition modulo \(\Lambda\). Equivalently: three points of \(E\), counted with
intersection multiplicity, are collinear if and only if their \(\Phi\)-preimages
in \(\C/\Lambda\) sum to \(0\).
\end{theorem}

The map \(\Phi\) is constructed in \cref{ch:embedding}; that it lands on \(E\)
comes from the differential equation for \(\wp\) in \cref{ch:weierstrass_wp};
that it is a bijection of Riemann surfaces is \cref{ch:embedding}; and the
collinearity statement, which is the heart of the matter, is proved in
\cref{ch:lines_divisors,ch:transport} using Abel's sum law from
\cref{ch:abel}.

\begin{keyidea}{Main conclusion}
If a line meets the cubic in three points \(P, Q, R\), counted with
intersection multiplicity, then those three points add to zero in the group law
transported from the torus:
\[
  P + Q + R = O .
\]
Solving for \(P + Q\) gives
\[
  P + Q = -R .
\]
Geometrically, this means: draw the line through \(P\) and \(Q\), take the
third intersection \(R\), then reflect \(R\) across the \(x\)-axis.
\end{keyidea}

\section{The Chain of Ideas}
\label{sec:intro-chain}

The proof of \cref{thm:main} is a chain of structural facts, each of which is
of independent interest and each of which occupies one part of this book. The
logical skeleton is shown in \cref{fig:dependency}.

\begin{figure}[ht]
\centering
\begin{tikzpicture}[node distance=1.2cm]
  \node[visualnode, text width=3.2cm] (torus) {complex torus\\\(T=\C/\Lambda\)};
  \node[visualnode, text width=3.3cm, right=1.4cm of torus] (wp) {functions\\\(\wp,\wp'\)};
  \node[visualnode, text width=3.2cm, right=1.4cm of wp] (cubic) {plane cubic\\\(E\subset\CP^2\)};
  \node[visualnode, text width=3.3cm, below=1.1cm of wp] (div) {principal divisors\\have Abel sum zero};
  \node[visualnode, text width=3.3cm, below=1.1cm of cubic] (line) {a line cuts\\\(P+Q+R\)};
  \draw[arrow] (torus) -- node[above, labelnode] {\(\Phi\)} (wp);
  \draw[arrow] (wp) -- (cubic);
  \draw[arrow] (line) -- (div);
  \draw[arrow] (div) -- node[left, labelnode] {\(P+Q+R=O\)} (torus);
  \draw[arrow] (cubic) -- (line);
\end{tikzpicture}
\caption{The chain of ideas behind \cref{thm:main}. The torus supplies the
group law; the Weierstrass functions \(\wp,\wp'\) carry it onto the cubic
\(E\); a line meeting \(E\) in \(P,Q,R\) records the divisor \(P+Q+R\); that
divisor is principal, so by Abel's sum law its points sum to \(0\) back on the
torus. The loop closes, and the chord--tangent law is forced.}
\label{fig:dependency}
\end{figure}

Read clockwise from the upper left, the chain runs as follows.

\begin{enumerate}
\item \textbf{The torus carries an obvious group law.} The quotient
\(\C/\Lambda\) inherits addition from \(\C\). Associativity, commutativity, and
inverses are immediate because they already hold in \(\C\)
(\cref{ch:lattices,ch:torus_group_surface}).

\item \textbf{The torus is a Riemann surface of genus one.} As a smooth compact
surface \(\C/\Lambda \cong S^1 \times S^1\) is a doughnut; as a Riemann surface
it has genus one. This is what eventually forces the cubic equation
(\cref{ch:torus_group_surface}).

\item \textbf{Meromorphic functions on the torus are scarce.} A nonconstant
holomorphic function cannot exist (Liouville), and the orders of zeros and poles
of a meromorphic function obey rigid global constraints
(\cref{ch:elliptic_liouville,ch:divisors}).

\item \textbf{Riemann--Roch counts those functions.} On a genus-one surface the
dimension counts are sharp enough to single out a function with a double pole
(\(\wp\)) and one with a triple pole (\(\wp'\)), and to force a single algebraic
relation between them (\cref{ch:riemann_roch,ch:weierstrass_wp}).

\item \textbf{That relation is the cubic.} The differential equation
\(\wp'^2 = 4\wp^3 - g_2\wp - g_3\) says precisely that \((\wp,\wp')\) traces out
\(E\). The map \(\Phi\) of \cref{thm:main} is the resulting embedding
(\cref{ch:weierstrass_wp,ch:embedding}).

\item \textbf{Lines cut principal divisors.} A line meeting \(E\) at \(P,Q,R\)
defines a meromorphic function on \(E\) whose divisor is
\(P + Q + R - 3O\). Pulled back to the torus, principal divisors have Abel sum
zero, which is exactly \(P + Q + R = O\)
(\cref{ch:abel,ch:lines_divisors,ch:transport}).
\end{enumerate}

Everything after that --- the explicit chord and tangent formulas, the role of
the point at infinity as the identity, the two-torsion points, the closed-form
addition laws --- is the working-out of this one loop
(\crefrange{ch:chord_tangent_cases}{ch:synthesis}).

\section{How to Read This Book}
\label{sec:intro-reading}

The book is organized into six parts that follow the chain above.

\begin{description}[leftmargin=1.4em, style=nextline]
\item[Part I --- Tori as Riemann Surfaces
  (\cref{ch:lattices,ch:torus_group_surface}).]
  Lattices, the quotient \(\C/\Lambda\), its complex structure, its group law,
  and its genus.
\item[Part II --- Elliptic Functions
  (\cref{ch:elliptic_liouville,ch:divisors}).]
  Doubly periodic meromorphic functions, Liouville's theorems, orders, and
  divisors.
\item[Part III --- Riemann--Roch on Genus One (\cref{ch:riemann_roch}).]
  The spaces \(L(D)\), the canonical divisor, and the dimension counts that
  make \(\wp\) and \(\wp'\) inevitable.
\item[Part IV --- The Weierstrass Functions
  (\cref{ch:weierstrass_wp,ch:embedding}).]
  Construction of \(\wp\), the differential equation, and the embedding
  \(\Phi\) into \(\CP^2\). The convergence, Eisenstein-series, and
  field-generation material of \cref{ch:weierstrass_wp} is the most technical in
  the book (graduate level) and may be skimmed on a first reading without losing
  the main thread.
\item[Part V --- The Group Law on the Cubic
  (\crefrange{ch:abel}{ch:chord_tangent_cases}).]
  Abel's sum law, lines as principal divisors, the transport of the group law,
  and the chord--tangent law case by case.
\item[Part VI --- Formulas, Computations, and Synthesis
  (\crefrange{ch:coordinate_formulas}{ch:synthesis}).]
  Closed-form addition and doubling formulas, the analytic addition formula, a
  fully worked numerical example, and pointers onward.
\end{description}

Four appendices support the main text.
\Cref{app:A} is a self-contained refresher on the complex analysis the book
relies on (Liouville's theorem, residues, the argument principle, Laurent
series); a reader who is rusty should skim it first.
\Cref{app:B} records how the present material grew out of a course's midterm and
final examinations, for readers who want that bridge.
\Cref{app:C} documents the Sage scripts that compute the exact coordinates for
every figure, so that all pictures are reproducible.
\Cref{app:D} contains full solutions to every exercise.

\paragraph{Prerequisites.}
We assume one course in single-variable complex analysis (holomorphic and
meromorphic functions, Laurent series, residues, Liouville's theorem, the
argument principle), comfort with quotient spaces from point-set topology, and
the basics of abelian groups and linear algebra. No prior exposure to Riemann
surfaces, divisors, or Riemann--Roch is assumed; these are developed here in the
single case of genus one.

\paragraph{Standing notation.}
The following symbols are fixed throughout. The complex numbers, reals, and
integers are \(\C\), \(\R\), \(\Z\); the upper half-plane is \(\HH\); and
\(\CP^n\) is complex projective space. A \emph{lattice} is written
\(\Lambda = \Z\omega_1 + \Z\omega_2\), and the associated \emph{complex torus}
is \(T = \C/\Lambda\), with \(O\) the image of \(0\), which will turn out to be
the group identity. The Weierstrass function and its derivative are \(\wp\) and
\(\wp'\); the invariants of \(\Lambda\) are \(g_2, g_3\), with discriminant
\(\Delta = g_2^3 - 27 g_3^2\). We write \(\Div\) for divisors, \(\divisor(f)\)
for the divisor of a function, \(\ord_p\) for order at \(p\), and \(L(D)\) for
the Riemann--Roch space of a divisor \(D\). Addition on the torus is written
\(+\); the chord--tangent law on \(E\) is written \(\oplus\) until
\cref{thm:main} identifies the two, after which we use \(+\) for both.

\section{Summary}
\label{sec:intro-summary}

The chord--tangent law looks like an arbitrary geometric construction, but it is
the image of ordinary addition on a complex torus under the Weierstrass
embedding \(\Phi\). \Cref{thm:main} is the precise statement, and
\cref{fig:dependency} is its proof in outline: the torus gives the group law,
\(\wp\) and \(\wp'\) carry it to the cubic, and the fact that a line cuts a
principal divisor closes the loop. The rest of the book builds each link of that
chain in turn.

\section*{Exercises}
\addcontentsline{toc}{section}{Exercises}

\begin{exercise}
\label{ex:intro-thirdpoint}
In the chord--tangent recipe, a line through two points \(P, Q\) of a cubic is
asserted to meet the curve in exactly one further point \(R\). Assuming the
points are counted with multiplicity and that we work in the projective plane,
explain informally why a third intersection point must exist. (The precise tool
is B\'ezout's theorem; see \cref{ch:lines_divisors}.)
\end{exercise}

\begin{exercise}
\label{ex:intro-tangent}
Explain why \emph{doubling} a point \(P\) (computing \(P \oplus P\)) uses the
tangent line at \(P\) rather than a chord. What happens to the chord through
\(P\) and \(Q\) as \(Q \to P\)?
\end{exercise}

\begin{exercise}
\label{ex:intro-liouville}
Recall Liouville's theorem: a bounded entire function is constant. Granting that
a meromorphic function on the torus \(\C/\Lambda\) with no poles is, in effect,
a bounded entire function, explain in one sentence why there is no nonconstant
holomorphic function on \(\C/\Lambda\). (This is proved carefully in
\cref{ch:elliptic_liouville}.)
\end{exercise}

\begin{exercise}
\label{ex:intro-map-arrows}
Using \cref{fig:dependency} and the chapter list in
\cref{sec:intro-reading}, match each of the five arrows in the diagram to the
chapter that establishes it.
\end{exercise}
